\section{Kernels on finite covers}\label{sec:kernels}

Let $p:Z\to C$ be a connected unramified principal $G$-cover, where
$G=D_n$ is the group of order $2n$ acting by the symmetries of a regular
$n$-gon, $n\geq3$. Write $V$ for its real reflection representation
and $\mathcal V=Z\times_G V$ for the associated flat Euclidean plane.
Rotation through $2\pi/n$ has no real invariant line, so $V$ is
irreducible. Moreover $\operatorname{End}_G(V)=\RR$. An endomorphism
commuting with this rotation is complex multiplication in the oriented
plane, and commutation with a reflection makes its multiplier real.

For any real Euclidean local system $\mathcal W$, denote by
$D^{\mathcal W}$ the twist of $D$ by $\mathcal W^*$. Its sections
are real homomorphisms $\mathcal W\to N$, and on a flat frame the
equation $D^{\mathcal W}S=0$ means that $D$ kills every column of $S$.
The Euclidean form identifies the dual with $\mathcal W$.
A transverse crossing on $\mathcal V$ at $t_0$ means that
$\ker D_{t_0}^{\mathcal V}$ is a real line and that differentiation
in $t$ induces an isomorphism from this line to the cokernel. Both
spaces have the same dimension by the index formula for an unramified twist.
For a map $u:X\to C$ from a connected curve, let $K_u\leq G$ be the
image of $\pi_1(X)$ under the monodromy of $u^*Z$. The choice of a
lift of the base point identifies parallel sections of $u^*\mathcal V$
with $V^{K_u}$.

\begin{theorem}\label{thm:kernels}
For a residual set of smooth paths of real Cauchy--Riemann operators
on $N$ with fixed symbol, every interior parameter where
$D_t^{\mathcal V}$ has a kernel is a transverse crossing. At any
such parameter let $S$ generate that kernel and put $D=D_t$. The map
$S:\mathcal V\to N$ is injective as a real map on every fiber. For
every smooth connected finite holomorphic cover $u:X\to C$, with
arbitrary branching,
\begin{equation}\label{eq:central-kernels}
 \ker u^*D
  =(u^*S)\bigl(H^0_{\rm flat}(X,u^*\mathcal V)\bigr)
  \simeq V^{K_u}.
\end{equation}
The path may have prescribed endpoints. Conversely, for every $p:Z\to C$
as above there exists a smooth path with a transverse crossing satisfying
these conclusions.
\end{theorem}

We first work at a prescribed section $S$ and later prove the generic
evaluation and the existence assertions. Suppose $D^{\mathcal V}S=0$
and that on a disk the two columns $s_1,s_2$ of $S$ in a flat frame
form a complex basis of $N$. Choose a closed subdisk in this disk with
nonempty interior. Let $\mathcal A_S$ be the space of smooth real
zeroth-order variations supported in the closed subdisk that annihilate
both columns. This closed real subspace of the space of smooth bundle
maps has its usual Fr\'echet topology.

\subsection{Perturbations preserving two solutions}
Write $F=\Lambda^{0,1}T^*C\otimes_{\CC}N$, and let $F^\dagger$
denote its real dual tensored with densities. The infinite-order Petri
condition at $x$ means that a real tensor of local solution and
adjoint-solution germs is zero whenever its pointwise product in
$N\otimes_{\RR}F^\dagger$ vanishes to infinite order at $x$.
Arbitrarily small unrestricted smooth real zeroth-order variations in
any neighborhood of $x$ attain this condition for the Cauchy--Riemann
symbol \cite[Definition~5.1 and Corollary~5.10]{wendl}. The formal-jet
version is given in \cite[Definition~1.6.3 and Theorems~1.6.17,~2.5.3]{dwbn}.

The condition at any point in the interior of the subdisk is dense
also in $D+\mathcal A_S$. To see this, start with an unrestricted
nearby operator $D'$ having the condition at $x$. On a small disk
about $x$ choose solutions $s'_j$ of $D's'_j=0$ close to $s_j$.
In a local frame $D'=\bar\partial+M'$ with $M'$ real linear.
If $T$ is the Cauchy transform on a disk of radius $r$, then
$\|T\|_{L^\infty\to L^\infty}\leq c r$. Thus
$w_j=-T(D's_j+M'w_j)$ has a solution for small $r$, and
$s'_j=s_j+w_j$ has the required equation. The sup norm of $w_j$
tends to zero as $D'$ approaches $D$. Interior elliptic estimates
on a smaller disk give closeness in any prescribed finite number of
smooth seminorms. The complex gauge $U$ sending $s_j$ to $s'_j$ is
then close to the identity. The conjugate $U^{-1}D'U$ has the same
symbol and kills $s_1,s_2$. The Petri condition is invariant under
this conjugation using the output action on $F$ and its contragredient
on $F^\dagger$. Multiply $U^{-1}D'U-D$ by a cutoff supported inside
the subdisk and equal to one near $x$. This proves the claimed density
at $D$ and, by the same construction, at every point of $D+\mathcal A_S$.

Consider an irreducible real orthogonal local system $\mathcal W$
of finite monodromy on $C$ minus finitely many points. Retain only
punctures with nontrivial inertia on $\mathcal W$.
At such a puncture use the orbifold disk whose stabilizer has the
order of the inertia action. This gives a Fredholm operator
$D^{\mathcal W}$ on orbifold sections. Equivalently, one can take
equivariant sections on the smooth compactification of the finite
normal cover defined by the monodromy. The two descriptions identify
kernels also after passage to a common normal cover. A bounded
invariant solution descends off the branch values and extends over
them by removability for the Cauchy--Riemann equation
\cite[Propositions~2.7.3 and~2.8.3]{dwbn}.

Let $w$ be a nonzero twisted solution with columns $w_1,\ldots,w_r$
in a real orthonormal flat frame of $\mathcal W$. If $\mathcal W$ is
inequivalent to $\mathcal V$, these columns together with $s_1,s_2$
are linearly independent as real global solutions on a common
connected normal cover. Indeed $w$ induces a nonzero intertwiner
from the irreducible fiber representation to the real solution space.
It is injective, and its image has zero intersection with the
inequivalent irreducible plane given by $S$. For $\mathcal W=\mathcal V$
the same independence holds for every nonzero $w$ orthogonal to $S$
in $L^2$. In this case two copies with a common nonzero vector
would coincide, and the two intertwiners with that image would differ
by a real scalar since $\operatorname{End}_G(V)=\RR$. In either case
strong unique continuation implies independence of the infinite
jets of these solutions at any regular point.

Choose a disk of flat frames in the variation region disjoint from
the branch values, and a point $x$ there at which $D$ satisfies the
Petri condition. We claim that the equation on nonzero sections,
linearized in sections and variations from $\mathcal A_S$, is surjective
at $w$. For $\mathcal W=\mathcal V$ use the $L^2$-complement of $S$
as section space. Omitting $S$ leaves the image of the section
derivative unchanged. An element of the annihilator consequently
solves the twisted adjoint equation. Write its columns on the disk
as $v_i$ in dual output densities. Testing against all local real
variations that kill $s_1,s_2$ gives the pointwise identity
\begin{equation}\label{eq:petri-constrained}
  \sum_{i=1}^r w_i\otimes v_i
      =s_1\otimes\alpha_1+s_2\otimes\alpha_2
\end{equation}
for smooth dual densities $\alpha_1,\alpha_2$. This is the annihilator
of the subspace of real maps vanishing on
$\operatorname{span}_{\RR}(s_1,s_2)$.
Both $\alpha_j$ satisfy the local adjoint equation. Indeed, test
\eqref{eq:petri-constrained} against $[D,\Phi]$ for a compactly
supported complex endomorphism $\Phi$ of $N$, with its induced action
on $F$. This commutator is zeroth order. Integration of the left
side gives $\sum_i\int\langle v_i,D(\Phi w_i)\rangle=0$.
The corresponding expression on the right is
$\sum_j\int\langle\alpha_j,D(\Phi s_j)\rangle$. The two test sections
$\Phi s_1$ and $\Phi s_2$ can be prescribed independently since
$s_1,s_2$ are a complex basis. Thus $D^\dagger\alpha_j=0$.
Taking infinite jets in \eqref{eq:petri-constrained}, the Petri
condition and independence of the input jets force every $v_i$
to vanish to infinite order. Adjoint unique continuation on the
connected punctured normal cover gives a zero annihilator. The
section derivative is Fredholm, so the asserted surjectivity follows.

\subsection{Exclusion of further kernels}
For each ramification type and representation, the distinct positions
of the effective punctures give local coordinates on its marked
configuration space. Write $s$ for their number and
$c_x=\dim_{\RR}(\mathcal W_x/\mathcal W_x^{I_x})$ for the codimension
of the invariants of local inertia on a nearby regular fiber.
The orbifold index formula \cite[Proposition~2.8.6]{dwbn} gives
\[
  \operatorname{ind}_{\RR}D^{\mathcal W}
   =(\dim_{\RR}\mathcal W)\operatorname{ind}_{\RR}D
       -(\operatorname{rank}_{\CC}N)\sum_x c_x
   =-2\sum_x c_x.
\]
In the incidence of projective kernel lines the moving punctures add
$2s$ parameters and passing to lines removes one. Thus the projection
index is $2s-1-2\sum_x c_x\leq-1$. For $\mathcal W=\mathcal V$ the
twist is unramified and fixed. Restriction to the complement of $S$
lowers its index by one more.

We spell out how to apply this count in the smooth constrained
parameter space. The finite monodromy groups, irreducible representations
up to equivalence, and homomorphisms from a punctured surface group
to a finite group give countably many discrete choices.
The compactified normal covers vary smoothly on a chart of distinct
branch positions by using the local maps $z\mapsto z^e$ at the
branches. Exhaust the configuration charts by countably many
compact subsets, chosen small enough that each subset misses some
open disk in the interior of the prescribed variation region. Here
it is the branch divisors that must miss the disk. The unramified
system $\mathcal V$ is a separate discrete choice.

Over each compact parameter subset the absence of a nonzero solution
(a complementary solution for $\mathcal V$) is open. Indeed
$L^2$-normalized solutions satisfy uniform elliptic estimates on
the smooth family of compact normal covers. To prove density, first
attain the Petri condition at a single point on the chosen open disk.
At that operator the normalized zeros over the compact parameter
subset form a compact set. By the surjectivity just proved, finitely
many allowed smooth operator variations suffice for surjectivity near
all these zeros. Shrinking their common parameter ball ensures that
all nearby zeros over the compact subset lie in those neighborhoods.
In section Sobolev completions the implicit function theorem now
gives finite-dimensional projective incidences over this ball,
of dimension strictly less than that of the ball by the index count.
Sard's theorem, applied on open chart neighborhoods, gives arbitrarily
small parameters with no zero over the compact subset. Local
normalization of kernel lines is legitimate since the scaling
direction maps to zero in the differential. Intersecting these open
dense conditions gives the asserted residual set in $D+\mathcal A_S$.

To deduce \eqref{eq:central-kernels}, pass to a finite normal cover
dominating $Z$ and the given $X$. Every irreducible summand in
its kernel supplies a twisted solution for the corresponding
monodromy system. Ineffective inertia can be omitted by descent
and removability. The only possible summand is $V$ pulled back
from $G$. Its multiplicity is one, since all such intertwiners
descend to $Z$ and the complement of $S$ in $\ker D^{\mathcal V}$
is zero. Descent through the sheet stabilizer of $X$ identifies
$\ker u^*D$ with the invariants of its image $K_u$ on this single
copy. Deleting preimages of the finite branch set does not affect
that image, since the lower cover $Z\to C$ is unramified.

\subsection{Generic paths and realization}
It remains to obtain the complex basis disk and fiberwise real
injectivity on generic paths. With $p$ fixed, consider the equation
$D^{\mathcal V}S=0$ at a nonzero $S$, varying $D$ without constraints
in this paragraph. On $Z$ two flat basis vectors map to independent
global real solutions by irreducibility. Their values are real
independent on an open dense subset. Indeed, on an open disk where two
values are proportional and one is nonzero, the ratio is a smooth
real function $a$. The product rule gives $(\bar\partial a)s_j=0$,
so $a$ is constant there, contradicting global independence by
unique continuation. Thus, on some disk away from a specified point,
variations of $D$ can prescribe the two inhomogeneous columns
independently. An annihilating adjoint vanishes on that disk and hence
everywhere. This proves universal transversality.

Evaluation of both columns at any marked point is submersive on
this incidence. To check the augmented equation, a distributional
adjoint annihilating the variations satisfies the homogeneous adjoint
equation off the mark and vanishes there by the same argument.
Its support is the mark, so it has finite order. Prescribe arbitrary
column values there for a test section and arrange that its image
under $D^{\mathcal V}$ vanishes to higher order than the distribution.
Such jets follow recursively from the solved $\bar\partial$ equation
in a coordinate, with the pure holomorphic Taylor coefficients free.
Cutoff Taylor polynomials realize them. Pairing with the augmented
equation first kills the evaluation multiplier and then, by elliptic
regularity, the point-supported adjoint.

For fixed $p$, the projection of the projectivized kernel incidence
to operators has index $-1$. Adding a path parameter therefore gives
regular isolated crossings on a generic path. A higher-dimensional
kernel would give a positive-dimensional projective fiber. At a
regular isolated crossing the time derivative maps the kernel line
isomorphically to the cokernel. The two evaluated real four-vectors
have rank one in real codimension three and rank zero in codimension
eight. Moving the marked point adds only two dimensions, so neither
stratum occurs at any crossing on a generic path. At one fixed point
the locus of complex rank one has real codimension two and the rank-zero
locus has codimension eight. Thus the columns also form a complex
basis on a disk at each crossing. All these rank strata are invariant
under nonzero real scaling, so the marked incidence calculation
passes to projective lines. Parametric transversality can be
applied with finitely many smooth variations on compact exhaustions
of the normalized solution charts as above, with cutoffs in the
interior of the time interval.

Mark a transverse crossing and choose a local linear normalization of
its generator. Variation of the section in a complement to its kernel
maps isomorphically onto the image of $D^{\mathcal V}$, and the time
variation spans the remaining cokernel. The implicit function theorem
therefore continues the marked crossing uniquely over an open path
chart. One may use a finite regularity completion and restrict to
smooth paths by elliptic regularity. On this chart, absence of the
additional twisted kernels over any one of the compact configuration
subsets is open. For the distinguished system the complement to the
continued line varies continuously. Density follows by perturbing
at the marked crossing while keeping its particular $S$ fixed, as
proved above, then multiplying the perturbation by a smooth time
cutoff equal to one near that crossing. The center and its generator
still solve the equation and hence are their own continuations.
Such crossing charts admit a countable cover, and their maps to
path space are local homeomorphisms. The exceptional subsets are
therefore meagre in path space. This proves the generic assertion
simultaneously at all interior crossings.

Finally start with an arbitrary connected principal $D_n$-cover
$Z\to C$. Choose a smooth real bundle map $S:\mathcal V\to N$.
Transversality to the real rank strata, of codimension three and
eight, gives a choice injective everywhere. At a fixed point we may
at the same time avoid the deficient complex rank strata. Choose any
real Cauchy--Riemann operator $D_*$. For a local flat $v$ prescribe
the real zeroth-order correction by
\[
                   P(Sv)=-D_*(Sv).
\]
Flat transition functions make this a smooth map on the subbundle
$S(\mathcal V)$, and a real bundle complement extends it to $N$.
Then $(D_*+P)^{\mathcal V}S=0$. The constrained perturbation proved
above gives an operator $D_0$ with this same $S$ and precisely the
kernels \eqref{eq:central-kernels}. Its twisted cokernel is a line.
On the complex basis disk an unrestricted real variation $Q$ can
prescribe both columns of $QS$ independently. A nonzero twisted
adjoint cannot vanish on the disk, so some such $Q$ pairs nontrivially
with it. The path $D_0+tQ$ has the required transverse crossing.
